% Source: https://artofproblemsolving.com/community/c3h2137372_equation_in_squares_given_sqrt_of_sum_wanted
% Stable identity: aops-topic-2137372-q01
Adding $8$ to both sides and factoring,
\[x^2y^2-4x^2-2y^2+8=\left(x^2-2\right)\left(y^2-4\right)=90 .\]
So $x^2-2$ is a divisor of $90$, and $x^2-2\ge-2$.
\begin{itemize}
\item $x=0$: $y^2-4=-45$, impossible. $x=\pm1$: $y^2-4=-90$, impossible.
\item $|x|\ge2$: then $x^2-2$ is a positive divisor of $90$ that is $2$ less than a perfect square. Of the positive divisors $1$, $2$, $3$, $5$, $6$, $9$, $10$, $15$, $18$, $30$, $45$, $90$, only $2$ works (as $2+2=4$). So $x=\pm2$, $y^2-4=45$ and $y=\pm7$.
\end{itemize}
Hence $x+y\in\{9,5,-5,-9\}$. For $\sqrt{x+y}$ to be real we need $x+y\ge0$, so the least possible value is $\boxed{\sqrt5}$ (from $x=-2$, $y=7$).
